\subsection{Multiplicidades en la descomposición primaria sobre un preesquema
algebraico.}
\label{subsection:IV.4.7-fr}

\begin{lemma}[4.7.1]
\label{IV.2.4.7.1-fr}
Sean $A$, $B$ dos anillos locales, $\mathfrak m$ el ideal maximal de $A$ y
$\rho:A\to B$ un homomorfismo local. Se supone que $B$ es un $A$-módulo plano
y que $B/\mathfrak mB$ es un $B$-módulo de longitud finita. Entonces, para que
un $A$-módulo $M$ sea de longitud finita, es necesario y suficiente que
$M_{(B)}$ sea un $B$-módulo de longitud finita, y se tiene
\begin{equation}
\operatorname{long}_B(M_{(B)})
=\operatorname{long}_A(M).\operatorname{long}_B(B/\mathfrak mB).
\tag{4.7.1.1}
\label{IV.2.4.7.1.1-fr}
\end{equation}
\end{lemma}

\begin{proof}
Es un caso particular de (2.5.5.2).
\end{proof}


\begin{corollary}[4.7.2]
\label{IV.2.4.7.2-fr}
Sean $A$, $B$ dos anillos noetherianos y $\rho:A\to B$ un homomorfismo de
anillos tal que $B$ es un $A$-módulo plano. Sean $\mathfrak q$ un ideal primo
minimal de $B$, $\mathfrak p=\rho^{-1}(\mathfrak q)$ y $M$ un $A$-módulo de
tipo finito~; entonces $\mathfrak p$ es un ideal primo minimal de $A$,
$M_{\mathfrak p}$ es un $A_{\mathfrak p}$-módulo de longitud finita y se tiene
\begin{equation}
\operatorname{long}_{B_{\mathfrak q}}(M_{(B)})_{\mathfrak q}
=(\operatorname{long}_{A_{\mathfrak p}}(M_{\mathfrak p}))
 (\operatorname{long}_{B_{\mathfrak q}}
 (B_{\mathfrak q}/\mathfrak pB_{\mathfrak q})).
\tag{4.7.2.1}
\label{IV.2.4.7.2.1-fr}
\end{equation}
\end{corollary}

\begin{proof}
En efecto, $B_{\mathfrak q}$ es un $A_{\mathfrak p}$-módulo plano
$(\mathbf{0}_{\mathrm I},6.3.2)$ y el homomorfismo
$A_{\mathfrak p}\to B_{\mathfrak q}$ es local, por lo que $B_{\mathfrak q}$ es
un $A_{\mathfrak p}$-módulo fielmente plano $(\mathbf{0}_{\mathrm I},6.6.2)$~; el hecho
de que $\mathfrak p$ es minimal en $A$ se deduce de (2.3.5). Además,
$A_{\mathfrak p}$ es entonces un anillo artiniano (Bourbaki, \emph{Alg. comm.},
chap.~IV, \S~2, n\textsuperscript{o}~5, prop.~9), y $M_{\mathfrak p}$ es, por
tanto, un $A_{\mathfrak p}$-módulo de longitud finita~; la fórmula (4.7.2.1) es,
pues, un caso particular de (4.7.1.1).
\end{proof}

\begin{proposition}[4.7.3]
\EGAReaderTarget{ega4.c04.prop.4-7-3}{IV.2.4.7.3-fr}
Sean $k$ un cuerpo de exponente característico $p$, $X$ un preesquema íntegro
localmente noetheriano y $K=\mathrm R(X)$ el cuerpo de funciones racionales de
$X$. Sean $k'$ una extensión de $k$, $(X'_\alpha)$ la familia de las
componentes irreducibles de $X'=X_{(k')}$ y $x'_\alpha$ el punto genérico de
$X'_\alpha$.
\begin{enumerate}
\item[\textup{(i)}] Se supone que $X'$ es localmente noetheriano \textup{(lo que ocurre si $X$
es \emph{localmente de tipo finito sobre} $k$, o si $k'$ es una extensión \emph{de tipo
finito} de $k$)}. Sean entonces $k''$ una extensión cualquiera de $k'$,
$X''=X\otimes_k k''$, $(X''_\beta)$ la familia de las componentes irreducibles
de $X''$ y $x''_\beta$ el punto genérico de $X''_\beta$. Si $x''_\beta$ está por
encima de $x'_\alpha$, se tiene
\begin{equation}
\operatorname{long}(\mathcal O_{x''_\beta})
=\operatorname{long}(\mathcal O_{x'_\alpha})
 .\operatorname{long}(\mathcal O_{Z'',x''_\beta})
\tag{4.7.3.1}
\label{IV.2.4.7.3.1-fr}
\end{equation}
donde se ha puesto $Z''=(X'_{\mathrm{red}})\otimes_{k'}k''$. En particular, si
$k''$ es una extensión \emph{separable} de $k'$ tal que $Z''$ es localmente
noetheriano, se tiene
$\operatorname{long}(\mathcal O_{x''_\beta})
=\operatorname{long}(\mathcal O_{x'_\alpha})$.

\oldpage[IV-2]{76}

\item[\textup{(ii)}] Si $X$ es localmente de tipo finito sobre $k$, o si $k'$ es una
extensión de tipo finito de $k$, los números
$\operatorname{long}(\mathcal O_{x'_\alpha})$ son potencias de $p$.

\item[\textup{(iii)}] Se supone que $X$ es localmente de tipo finito sobre $k$.
Entonces, si $k'$ es perfecto, todos los números
$\operatorname{long}(\mathcal O_{x'_\alpha})$ son iguales y solo dependen de la
extensión $K$ de $k$ \textup{(y no de la extensión perfecta particular $k'$ de $k$
considerada)}. Además, existe una extensión radicial finita $k_1$ de $k$ tal
que, si $x_1$ es el punto genérico de $X_1=X\otimes_k k_1$ (que es
irreducible), se tiene
$\operatorname{long}(\mathcal O_{x'_\alpha})
=\operatorname{long}(\mathcal O_{x_1})$ para todo $\alpha$.
\end{enumerate}
\end{proposition}

\begin{proof}
\begin{enumerate}
\item[(i)] Recordemos (4.2.4) que las componentes irreducibles $X'_\alpha$,
que son en número finito, corresponden biunívocamente a los ideales primos
minimales $\mathfrak p_\alpha$ de $K\otimes_k k'$, y que
$\mathcal O_{x'_\alpha}$ es un anillo \emph{artiniano}, isomorfo a
$(K\otimes_k k')_{\mathfrak p_\alpha}$~; estos anillos solo dependen, por tanto,
de las extensiones $K$ y $k'$ de $k$ (y son, en consecuencia, los mismos
($k'$ fijado) para todos los $k$-preesquemas íntegros que tengan el mismo cuerpo
$K$ de funciones racionales). La fórmula (4.7.3.1) se deduce de (4.7.2.1)
aplicada al anillo noetheriano $A$ de una vecindad afín de $x'_\alpha$ en $X'$, y
al anillo $B$ de una vecindad afín suficientemente pequeña de $x''_\beta$ en
$X''$, siendo $\mathfrak p$ y $\mathfrak q$ los ideales primos minimales
correspondientes a $x'_\alpha$ y $x''_\beta$, respectivamente, y tomando $M=A$~:
se observará que $A/\mathfrak p$ es el anillo de una vecindad afín de $x'_\alpha$ en
$X'_{\mathrm{red}}$, y que $B=A\otimes_{k'}k''$~; si
$\mathfrak q'=\mathfrak q/\mathfrak pB$ es el ideal primo minimal de
$B/\mathfrak pB$, se tiene
$\mathcal O_{Z'',x''_\beta}
=(B/\mathfrak pB)_{\mathfrak q'}
=B_{\mathfrak q}/\mathfrak pB_{\mathfrak q}$,
de donde se obtiene (4.7.3.1). La última afirmación de (i) se deduce de que, en
este caso, $Z''$ es reducido (4.2.4) y, por tanto,
$\mathfrak pB_{\mathfrak q}=\mathfrak qB_{\mathfrak q}$.

\item[(ii)] Si $k'$ es una extensión de tipo finito de $k$, existe una
extensión puramente trascendente $k'_0$ de $k$ tal que $k'$ es una extensión
finita de $k'_0$~; como $k'_0$ es separable sobre $k$, se deduce de (i),
aplicado reemplazando $k'$ y $k''$ por $k$ y $k'_0$, respectivamente, que se
puede reemplazar $X$ por $X\otimes_k k'_0$; dicho de otro modo, se puede suponer
que $k'$ es una extensión \emph{finita} de $k$. Existe entonces una extensión finita
cuasigaloisiana $k''$ de $k$ que contiene a $k'$, y la fórmula (4.7.3.1)
muestra que basta demostrar la afirmación para $k''$. Pero $k''$ es una
extensión \emph{separable} de una extensión radicial finita $k_1$ de $k$, por lo que,
en virtud de (i), la cuestión se reduce en este caso a demostrar la afirmación
cuando $k'=k_1$ es además \emph{radicial}. Se sabe entonces que $K\otimes_k k_1$ es un
anillo artiniano que tiene un solo ideal primo minimal (4.3.2)~; si
$\overline K$ es una clausura algebraica de $K$, la longitud de
$K\otimes_k k_1$ divide la del anillo artiniano $\overline K\otimes_k k_1$ en
virtud de (4.7.1.1). Pero $\overline K\otimes_k k_1$ es una $\overline K$-álgebra
artiniana que tiene un solo ideal primo (necesariamente maximal), y su cociente
por este ideal es una extensión finita de $\overline K$ y, por tanto,
necesariamente igual a $\overline K$~; la longitud de
$\overline K\otimes_k k_1$ es, pues, igual a su rango sobre $\overline K$, es
decir, a $[k_1:k]$, que es una potencia de $p$.

Supongamos a continuación que $X$ es localmente de tipo finito sobre $k$~;
entonces $X''$ es localmente noetheriano para toda extensión $k''$ de $k'$.
Tomando como $k''$ una extensión algebraicamente cerrada de $k'$, la cuestión se reduce, en virtud de (4.7.3.1), a demostrar la afirmación cuando $k'$ es
\emph{algebraicamente cerrado}. Como $k'$ es entonces una extensión \emph{separable} de la
clausura algebraica $\overline k$ de $k$, se deduce nuevamente de (i) que
la cuestión se reduce al caso en que $k'=\overline k$. Ahora bien, se sabe
(reemplazando $X$ por un abierto afín de tipo finito sobre $k$) que existe una
extensión radicial

\oldpage[IV-2]{77}

finita $k_1$ de $k$ tal que $(X\otimes_k k_1)_{\mathrm{red}}$ es separable sobre
$k_1$ (4.6.6)~; por tanto,
$(X\otimes_k k_1)_{\mathrm{red}}\otimes_{k_1}\overline k$ es reducido (4.6.4).
En virtud de (4.7.3.1), la cuestión se reduce de nuevo al caso en que $k'$ es una
extensión finita radicial de $k$, y se concluye como antes.

\item[(iii)] Observemos que dos extensiones $k'$, $k''$ de $k$ están contenidas
en una misma extensión algebraicamente cerrada de $k$~; para demostrar que el
conjunto de los números $\operatorname{long}(\mathcal O_{x'_\alpha})$, para
$k'$ \emph{perfecta}, no depende de la elección de $k'$, es posible limitarse, por
tanto, a comparar estos números para $k'$ y para una extensión algebraicamente
cerrada $k''$ de $k'$~; en este caso, la afirmación se deduce de (i), puesto que
$k''$ es \emph{separable} sobre $k'$. Para demostrar que todos los números
$\operatorname{long}(\mathcal O_{x'_\alpha})$ son iguales, es posible limitarse,
por tanto, al caso en que $k'=\overline k$, clausura algebraica de $k$. Ahora
bien, determinada entonces la extensión radicial finita $k_1$ de $k$ como en
(ii), los $x'_\alpha$ están todos por encima del único punto genérico de
$X_1=X\otimes_k k_1$, y la relación
$\operatorname{long}(\mathcal O_{x'_\alpha})=
\operatorname{long}(\mathcal O_{x_1})$ para todo $\alpha$ se deduce nuevamente
de (4.7.3.1), dada la elección de $k_1$.
\end{enumerate}
\end{proof}
